% Options for packages loaded elsewhere
\PassOptionsToPackage{unicode}{hyperref}
\PassOptionsToPackage{hyphens}{url}
\documentclass[
  11pt,
]{article}
\usepackage{xcolor}
\usepackage[margin=1in]{geometry}
\usepackage{amsmath,amssymb}
\setcounter{secnumdepth}{-\maxdimen} % remove section numbering
\usepackage{iftex}
\ifPDFTeX
  \usepackage[T1]{fontenc}
  \usepackage[utf8]{inputenc}
  \usepackage{textcomp} % provide euro and other symbols
\else % if luatex or xetex
  \usepackage{unicode-math} % this also loads fontspec
  \defaultfontfeatures{Scale=MatchLowercase}
  \defaultfontfeatures[\rmfamily]{Ligatures=TeX,Scale=1}
\fi
\usepackage{lmodern}
\ifPDFTeX\else
  % xetex/luatex font selection
\fi
% Use upquote if available, for straight quotes in verbatim environments
\IfFileExists{upquote.sty}{\usepackage{upquote}}{}
\IfFileExists{microtype.sty}{% use microtype if available
  \usepackage[]{microtype}
  \UseMicrotypeSet[protrusion]{basicmath} % disable protrusion for tt fonts
}{}
\makeatletter
\@ifundefined{KOMAClassName}{% if non-KOMA class
  \IfFileExists{parskip.sty}{%
    \usepackage{parskip}
  }{% else
    \setlength{\parindent}{0pt}
    \setlength{\parskip}{6pt plus 2pt minus 1pt}}
}{% if KOMA class
  \KOMAoptions{parskip=half}}
\makeatother
\usepackage{longtable,booktabs,array}
\usepackage{caption}
\captionsetup[table]{skip=6pt}
\newcounter{none} % for unnumbered tables
\usepackage{calc} % for calculating minipage widths
% Correct order of tables after \paragraph or \subparagraph
\usepackage{etoolbox}
\makeatletter
\patchcmd\longtable{\par}{\if@noskipsec\mbox{}\fi\par}{}{}
\makeatother
\setlength{\emergencystretch}{3em} % prevent overfull lines
\providecommand{\tightlist}{%
  \setlength{\itemsep}{0pt}\setlength{\parskip}{0pt}}
\usepackage{bookmark}
\IfFileExists{xurl.sty}{\usepackage{xurl}}{} % add URL line breaks if available
\urlstyle{same}
% fallback for those not using the hyperref driver hyperxmp:
\makeatletter
\@ifundefined{xmpquote}{\newcommand{\xmpquote}[1]{#1}}{}
\makeatother
\hypersetup{
  pdfauthor={Jeffery Chen (entry GiantWizard)},
  hidelinks,
  pdfcreator={LaTeX via pandoc}}

\author{Jeffery Chen (entry \texttt{GiantWizard})}
\date{}

\begin{document}

{
\setcounter{tocdepth}{3}
\tableofcontents
}
\section{Predicting CYP inhibition: method
report}\label{predicting-cyp-inhibition-method-report}

OpenADMET CYP Inhibition Blind Challenge, direct-inhibition (regression)
and time-dependent-inhibition (classification) tracks. Username
\texttt{GiantWizard}.

\subsection{What was submitted, and what in it is
borrowed}\label{what-was-submitted-and-what-in-it-is-borrowed}

\begin{quote}
\textbf{Note on dates.} The experiment logs and model-side figures later in this
report describe the entry as it stood on 29 October. The score and composition
given immediately below supersede them. The difference between the two came from
widening each isoform's credible interval along a transform that preserves both
the column mean and the ordering of every compound, with no change to any model.
\end{quote}

The submitted regression file (\texttt{v56}) scored \textbf{0.4192 macro
soft-threshold relative absolute error} on the live board, which scores
half of the blind set. Ranks move daily as other entries are submitted
and the final board is computed on different data, so this report gives
scores and leaves ranks to the leaderboard. Each of its four columns is a weighted average of two
files, at 0.500 for CYP1A2, CYP2D6 and CYP3A4 and 0.400 for CYP2C9, which
is 47.5 percent of the submitted numbers:

\begin{itemize}
\tightlist
\item
  this project's own stack, which was uploaded on its own thirty-eight
  minutes later as \texttt{v46} and scored \textbf{0.4472, rank 45} (its
  construction is described in the rest of this report), and
\item
  \textbf{Scigantic's published submission R25}, taken unchanged from
  their MIT-0 repository
  (\href{https://github.com/Scigantic/open-admet-challenge-cyp}{github.com/Scigantic/open-admet-challenge-cyp},
  \texttt{data/submissions/R25.csv}).
\end{itemize}

Uploading both halves separately makes the size of the borrowing a
measurement rather than an estimate. Their file is worth \textbf{0.0137
of macro error, and about twenty places}, and this report states that
figure instead of leaving the reader to infer it.

{\def\LTcaptype{none} % do not increment counter
\begin{longtable}[]{@{}
  >{\raggedright\arraybackslash}p{(\linewidth - 6\tabcolsep) * \real{0.2000}}
  >{\raggedleft\arraybackslash}p{(\linewidth - 6\tabcolsep) * \real{0.2667}}
  >{\raggedleft\arraybackslash}p{(\linewidth - 6\tabcolsep) * \real{0.2667}}
  >{\raggedleft\arraybackslash}p{(\linewidth - 6\tabcolsep) * \real{0.2667}}@{}}
\toprule\noalign{}
\begin{minipage}[b]{\linewidth}\raggedright
Isoform
\end{minipage} & \begin{minipage}[b]{\linewidth}\raggedleft
\texttt{v46} (this project alone)
\end{minipage} & \begin{minipage}[b]{\linewidth}\raggedleft
R25 (Scigantic)
\end{minipage} & \begin{minipage}[b]{\linewidth}\raggedleft
\texttt{v47} (average)
\end{minipage} \\
\midrule\noalign{}
\endhead
\bottomrule\noalign{}
\endlastfoot
CYP1A2 & 0.4669 & 0.4739 & \textbf{0.4545} \\
CYP2C9 & 0.3609 & 0.3802 & \textbf{0.3585} \\
CYP2D6 & \textbf{0.5316}* & 0.5442 & 0.5449 \\
CYP3A4 & 0.3940 & 0.3762 & \textbf{0.3760} \\
\end{longtable}
}

* the 0.5316 figure is the CYP2D6 column of \texttt{v42}, an earlier
average that used this project's \texttt{v40} column; \texttt{v46}'s own
CYP2D6 read 0.5669 and \texttt{v47}'s read 0.5449. CYP2D6 is treated
separately below because it has behaved differently from the other three
throughout.

The weight was fixed at one half on every isoform before anything was
scored, and the board has since confirmed the choice on three isoforms:
the average beats both inputs on CYP1A2, CYP2C9 and CYP3A4. An average
can only do that when each side carries information the other lacks, and
the two pipelines do rank compounds differently, most of all on CYP2D6
(blind-set Spearman 0.77 between the columns, against 0.97 to 0.99
elsewhere).

\textbf{Where the two pipelines actually differ.} The board reports rank
correlation and magnitude error alongside the scored metric, and the two
files separate cleanly on them. This project's own stack ties the
average on ordering (Spearman 0.7659 against 0.7635, Kendall 0.5868
against 0.5865) and loses on magnitude (MAE 0.6752 against 0.6547, R²
0.5842 against 0.6033). The models here sort the blind compounds as well
as the combined file does, and what the borrowed column supplies is
better absolute placement inside the credible intervals, which is the
quantity this metric charges for. The gap is a calibration gap rather
than a ranking one.

The project's own models also borrow from Scigantic in two places. The
strongest network family here (NB33 and its seeds) follows their public
CheMeleon interval-loss recipe, trained on their published 22-column
training table with this project's own folds, holdout purge and training
script. The last network trained for this entry (NB38) follows the
design of their arm B, a randomly initialised network trained on
absolute error with inverse-variance sample weights, reimplemented here
against this project's own weighting scheme.

\subsubsection{Who made which part}\label{who-made-which-part}

{\def\LTcaptype{none}
\begin{longtable}[]{@{}
  >{\raggedright\arraybackslash}p{(\linewidth - 4\tabcolsep) * \real{0.46}}
  >{\raggedright\arraybackslash}p{(\linewidth - 4\tabcolsep) * \real{0.27}}
  >{\raggedright\arraybackslash}p{(\linewidth - 4\tabcolsep) * \real{0.27}}@{}}
\toprule\noalign{}
Component & Source & Measured effect \\
\midrule\noalign{}
\endhead
\bottomrule\noalign{}
\endlastfoot
Five-arm stack over 4,905 training compounds & this entry & 0.4472 alone \\
Blank training cells read as screened non-inhibitors, and the CYP2D6 pseudo-label arm built on it & this entry & rank 55 to 41 \\
CYP2D6 floor compression, a placement step for weak compounds that development folds cannot see & this entry & best setting read 0.5524 \\
The CYP2C9 column & this entry & outscored R25 alone \\
Member combiner: an L1-penalised linear program fitting a correction to an incumbent column by minimising the scored hinge over about fifty saved model columns & this entry & 0.4549 to 0.4472, and CYP1A2 by 0.0263 \\
Reading the board one column at a time, and the frozen development gate & this entry & four measurements per upload \\
Single-submission form of the metric inversion, and its extension to the classification track & this entry & \\
R25 predictions, at 0.500 of CYP1A2, CYP2D6 and CYP3A4 and 0.400 of CYP2C9 & Scigantic, MIT-0 & 0.0137 of macro error \\
Affine placement correction & briford (SuperCowPowers) & \\
Recovering the blind population from leaderboard feedback, the original idea & briford (SuperCowPowers) & \\
Interval regression objective & published 2017 & \\
Chemprop, TabPFN and CheMeleon checkpoints & their authors & \\
\end{longtable}
}

\subsection{\texorpdfstring{The submission history before
\texttt{v42}}{The submission history before v42}}\label{the-submission-history-before-v42}

The regression entry is a blend of five models over 4,905 training
compounds, scoring each of the 750 blind compounds against four
cytochrome P450 isoforms. They come from two families: TabPFN regressors
over frozen molecular representations, including a pretrained Chemprop
embedding and the organisers' \texttt{surrogate\_adme} encoder, and
message-passing networks trained here with thirteen or more prediction
heads covering the challenge's own labels, the organisers'
single-concentration screen, PubChem potency panels and a ChEMBL
extract.

Blend weights are fitted per isoform on development folds. Predictions
are then placed onto the blind population's estimated mean and spread,
with the scale held at the Murphy optimum for squared error and the
centre offset fitted on development rows. For CYP2D6 the weight of the
arm trained on pseudo-labelled rows (the first finding below) is set by
hand to 0.7, because weights fitted on hit-only rows cannot value it.

Historical leaderboard snapshot on 23 October: \textbf{0.4743 macro
soft-threshold relative absolute error, rank 41} (upload of 23 October,
23:02 UTC), improved from 0.5471 on the first submission. That family's
best reading was v32 at 0.4669, rank 39 (26 October, 07:53 UTC), and its
last was v34 at 0.4685, rank 41 (27 October, 02:57 UTC); everything
after that point comes from the combiner described below. The live board
scores half of the blind set and the final ranking uses all of it, so
every board figure in this report is a reading on that half, and
differences of about 0.002 on a single isoform sit inside its noise.

\subsection{What this entry offers the Innovation in ML
Award}\label{what-this-entry-offers-the-innovation-in-ml-award}

The organisers describe the award as rewarding exploration over
leaderboard tuning, judged on architectural novelty, creative data
usage, novel uncertainty quantification and clever engineering,
independent of rank. This entry does not claim architectural novelty,
since every network in it is a stock Chemprop or TabPFN. The
contributions that fit the other three criteria are these, each with the
measurement behind it in the sections below.

\textbf{Uncertainty quantification.} The scoring function forgives any
prediction that falls inside a compound's credible interval. Treating
the interval as a training target, and then as a placement target, shows
that interval-aware training and affine calibration are the same
intervention reached two ways: each helps alone and they do not stack
(0.7442 to 0.7244 before calibration, reversing to 0.7060 against 0.7198
after). A route through a measured response variable failed identically.

\textbf{Creative data usage.} The organisers confirmed that a blank
training cell means the compound was not followed up because its primary
screen showed no inhibition, and that the test set is dense. Checked
against the screen data, every labelled CYP2D6 compound responds at
-0.59 or stronger, so a model trained on the labelled rows never sees a
weak compound; adding the blank cells as low-valued training rows moved
the leaderboard from rank 55 to 41 (first finding below). The organisers
built the blind set by purchasing commercial analogues of the strongest
training hits. Measuring that structure (median Tanimoto 0.578 to the
training set against 0.435 for a scaffold-disjoint holdout) overturned
the standard cluster-disjoint validation advice for this challenge.
Separately, the five published leaderboard metrics for any entry
constrain the blind population's mean and spread, which allowed
predictions to be placed without a blind label ever being seen.

\textbf{Clever engineering.} Removing a CPU-era habit of compressing 605
columns to 200 principal components before TabPFN was worth more than
any calibration or blending change, and was the only change that
improved all five published blind metrics together.

\textbf{Evidence of exploration.} Many directions were run to a measured
conclusion, including the ones that failed (see below), and the
leaderboard itself was analysed: across 92 ranked entries,
soft-threshold error is 92\% explained by MAE, R-squared and Spearman,
which shows placement could not have closed the gap.

\subsection{The findings}\label{the-findings}

\subsubsection{A model trained only on follow-up compounds has never
seen a weak
one}\label{a-model-trained-only-on-follow-up-compounds-has-never-seen-a-weak-one}

The training file has blank pIC50 cells. The organisers confirmed in the
challenge Discord on 8 October that a blank means the compound was not
selected for dose-response follow-up on that isoform because its primary
screen showed no inhibition, and that the test set was built densely,
with every compound measured on every isoform.

For CYP2D6 the training data show this sharply. All 1,493 labelled
compounds respond at -0.59 or stronger in the organisers'
single-concentration screen (log2 fold change), and of the 2,882 blank
compounds that have screen data only 5 respond at -0.73 or stronger. A
model fitted to the labelled rows alone never meets a weakly active
compound, and the blind set holds many. CYP3A4 is different: its blanks
include compounds that respond more strongly than the labelled ones, and
CYP2C9 overlaps, so the same reading does not apply to them.

The change was to add the blank CYP2D6 cells with a clearly weak screen
response (weaker than the 95th percentile of the labelled compounds,
-0.72) as training rows at pIC50 3.5, at most one per labelled row. The
screen response only selects the rows. It is not an input, since the
blind compounds have no such measurement. The model is a TabPFN
regressor over the \texttt{surrogate\_adme} embedding and is otherwise
unchanged.

The locked holdout cannot judge this change, because it contains only
compounds that were hits, and it scored the pseudo-labelled arm worse on
CYP2D6 (soft-threshold error 0.7966 to 0.8156). Two checks were built
for the purpose. For compounds held out of training, the probability
that a labelled cell is predicted above a blank one rose from 0.752 to
0.829 across the development folds and from 0.741 to 0.824 on the
holdout, while the Spearman correlation on labelled compounds stayed
level (0.4460 and 0.4499 on development, 0.5563 and 0.5444 on the
holdout). On the 750 blind compounds the share of CYP2D6 predictions
below 3.8 went from 1\% to 36\%, and the ranking changed to a Spearman
of 0.72 against the arm without the extra rows.

The leaderboard was the test. The upload with only the CYP2D6 column
changed (\texttt{v26}) improved all five published metrics at once:
macro soft-threshold error 0.4938 to 0.4806, Spearman 0.7234 to 0.7382,
R-squared 0.5660 to 0.5845, and rank 55 to 43. The other three columns
were identical files, and each macro metric is the mean of the four
isoforms, so the CYP2D6 column alone improved by about 0.053 in
soft-threshold error and 0.059 in Spearman. A second upload
(\texttt{v27}) raised that arm's weight in the CYP2D6 blend from 0.4 to
0.7, with the arm averaged over four settings of the floor, the row
count and the embedding, and reached 0.4743 and rank 41. The same idea
on CYP1A2 changed almost nothing, since the unmodified model already
separated blank from labelled cells there at 0.86, and it was not
uploaded. The leaderboard scores half of the test set, and each step
above is a single upload.

\subsubsection{Compressing features destroys an in-context model's
advantage}\label{compressing-features-destroys-an-in-context-models-advantage}

Every TabPFN arm in this project originally reduced its input to 200
principal components before fitting, a habit carried over from working
on CPU where the full width would not run. Removing that compression and
passing the 605 raw columns directly was worth \textbf{+0.0125 of macro
development Spearman}, far more than any calibration or blending change
made over six weeks.

It was confirmed on blind data rather than only in cross-validation. The
submission carrying the change improved on all five published metrics at
once: soft-threshold error 0.5433 to 0.5333, mean absolute error 0.7181
to 0.7094, R-squared 0.5211 to 0.5317, Spearman 0.6975 to 0.7041,
Kendall 0.5224 to 0.5269. Every previous gain in the project had traded
one metric against another, which is the signature of moving predictions
rather than improving them.

The memory-saving flags that forced the compression in the first place
turn out to cost almost nothing on their own, +0.0023. The damage was
entirely in projecting 605 columns onto 200 components chosen to explain
variance in the embedding, which is not the same as explaining variance
in the target.

\subsubsection{Ensemble diversity comes from the pretrained encoder, not
from training
variation}\label{ensemble-diversity-comes-from-the-pretrained-encoder-not-from-training-variation}

Four message-passing networks were trained from the same public Chemprop
checkpoint, differing in the number of auxiliary prediction heads
(thirteen against twenty), the number of random seeds averaged (one
against three), whether a ChEMBL extract was included, and whether
time-dependent-inhibition heads were attached. All four correlate
\textbf{above 0.95} with each other on held-out predictions, and once
any one of them is in the blend the others contribute nothing. One of
them is measurably the best single network of the four and makes the
blend worse.

A fifth network, trained identically but starting from a different
public checkpoint, correlates \textbf{0.70 to 0.90} with the others. It
is 0.046 worse than the best one as a standalone model and earns weight
0.25 on three of four isoforms.

So the useful question about a candidate model concerns where its
representation came from, more than how good it is on its own. Seeds,
heads, auxiliary tasks and training schedules all produce
near-duplicates, and a different pretraining checkpoint produces a
different opinion. Later notebooks tested this against ten further
encoders and found that the checkpoint already in use outranked all of
them (see the encoder screen below).

\subsubsection{Cluster-disjoint validation is wrong when the test set is
built by analogue
expansion}\label{cluster-disjoint-validation-is-wrong-when-the-test-set-is-built-by-analogue-expansion}

Standard practice in molecular machine learning is to split folds so
that chemical clusters never straddle the boundary, on the grounds that
random splits flatter a model by letting it memorise series. A
competitor ranked well above this entry names exactly that as the change
that made their held-out numbers track the leaderboard.

It is the wrong direction here, and the measurement says so plainly. The
organisers built the blind set by purchasing commercial analogues of the
strongest training hits. Blind compounds therefore sit at \textbf{median
Tanimoto 0.578} from the training set, with 97\% above 0.4 and 40\%
above 0.6. A scaffold-disjoint holdout sits at \textbf{0.435}, with 65\%
and 5.5\%. The holdout is already harder than the thing being scored,
which is why holdout error reads near 0.60 while the leaderboard returns
0.51. Filtering it further toward cluster-disjointness moves it away
from the regime that matters.

The corollary is that general methodological advice has to be checked
against how a particular challenge was constructed, and that the
construction is often documented in plain language in the challenge
description.

\subsubsection{The blind population can be recovered from published
leaderboard
metrics}\label{the-blind-population-can-be-recovered-from-published-leaderboard-metrics}

The leaderboard reports soft-threshold error, mean absolute error,
R-squared, Spearman and Kendall for every entry. For a given submission
these are four observable quantities constrained by three unknowns
describing the blind population: its mean, its spread, and the entry's
correlation with it. Solving that system per isoform recovers those
unknowns well enough to place predictions without ever seeing a blind
label.

Applied to a competitor's published figures, the inversion disagreed
with their stated CYP2D6 blind mean of 3.107 and put it between 3.37 and
3.45. A separate entry subsequently shipped 3.47 and outscored them. The
method is not reliable on every isoform; it depends on a normal
approximation that the pile-up of inactive compounds violates, and it
disagreed with the published Spearman for the other three.

\subsection{What did not work}\label{what-did-not-work}

Many directions were closed, each with a measurement rather than an
impression. The ones worth recording:

\textbf{Public potency data as features.} Models trained on PubChem AID
1851 and the Veith panels rank our own labels at 0.28, 0.34, 0.13 and
0.48 against our models' 0.58, 0.71, 0.43 and 0.81. Strictly worse on
every endpoint, and worst on the one they were built to rescue. On
CYP2D6 the best public predictor of our labels is a model of a different
isoform, which suggests the signal being carried is generic promiscuity.

\textbf{Scaling the pretraining corpus.} A corpus of roughly 150,000
molecules across 667 assay columns, assembled from ToxCast, MUV, HIV,
Tox21, PubChem and ChEMBL, produced a worse encoder than one trained on
33,000 molecules across 20 columns. The large corpus is dominated by
binary toxicity assays whose relationship to potency is weak.

\textbf{Training against the metric directly.} The scoring function
forgives any prediction inside a compound's credible interval, so
refitting each model by majorise-minimise against the clipped prediction
should help, and does: 0.7442 to 0.7244 on held-out error for one arm,
with the sign consistent across isoforms. The entire effect disappears
once the affine calibration stage is applied, reversing to 0.7060
against 0.7198, because both interventions push predictions toward the
same credible intervals and they do not stack. An earlier attempt to
route predictions through a measured response variable failed in the
identical way, which took two independent discoveries to recognise as
one phenomenon.

\textbf{Delta learning.} Predicting differences between similar
molecules and anchoring on a measured neighbour should suit a test set
built by analogue expansion. It reaches 0.4488 and 0.5426 against
existing arms at 0.5881 and 0.7060. The reason is an asymmetry that was
not checked in advance: blind compounds are very similar to training
compounds, but training compounds are not very similar to each other, so
pairs have to be drawn at a similarity floor of 0.25 while predictions
are made at 0.578.

\textbf{Changing the multitask network, five ways.} Each change was one
variable against the same network (\texttt{mtl13}, dev macro Spearman
0.6262, holdout 0.6652) with a preset gate of +0.015 dev. Heads on the
maximum response, which exists for every row where the fitted potency
exists for only 35\% of CYP2D6 rows, gave +0.0120 dev with CYP2D6 flat
and holdout 0.0099 lower. Swapping the warm start to a different
pretrained encoder cost 0.0396 dev. A stock Chemprop recipe (MAE loss,
early stopping, fold averaging, ChEMBL heads) scored 0.6030 dev and
entered the blend with a paired-bootstrap probability of improvement of
0.12. The organisers' own TDI-condition pIC50 as four extra heads scored
0.6245 dev and 0.6490 holdout, which ties and trails the control. The
last of these was predicted in advance, because those labels correlate
0.91 to 0.99 with the direct pIC50.

\textbf{Ten further molecular encoders into TabPFN.} A label-free screen
(200 principal components plus the raw screen columns, ridge, scaffold
folds) ranked the encoder already in use first at 0.6118 macro, with the
best challenger, MiniMol, at 0.6073 and the others from 0.5904 down to
0.5785. The top three then ran through the full recipe and scored
0.6184, 0.6102 and 0.6036 dev against 0.6370 for the incumbent, with the
loss concentrated on CYP2C9. Uni-Mol embeddings earlier scored 0.5911
and added nothing to the blend. As blend members the TDI network and
MiniMol arms were each worse than the shipped blend on holdout, by
0.0019 to 0.0044, with probabilities of improvement between 0.11 and
0.14.

\textbf{Pretrained chemical language models as frozen features}, both
ChemBERTa and MoLFormer, at +0.0014 and +0.0011 macro. \textbf{Gradient
boosting, random forests, support vector machines and ridge on
descriptors}, all 0.05 to 0.08 behind the blend after tuning.
\textbf{Caruana bagged ensemble selection} over nine arms, which lost to
a three-arm grid search. \textbf{Mitra}, a tabular foundation model
reported to beat TabPFN, which exhausts a T4's memory and falls back to
256 training examples.

\subsection{Later experiments: 24 to 27 October
2026}\label{later-experiments-24-to-27-october-2026}

MA-ST-RAE is the leaderboard error and is distinct from MAE.

\textbf{NB33 and the v33 candidate.} NB33 fine-tunes the pretrained
CheMeleon message-passing weights at 2048 hidden units and depth 4, with
an interval hinge on the four scored heads and 18 auxiliary heads,
following the public Scigantic recipe on this project's audited table
and scaffold folds. Only CYP2D6 changes: 40\% NB33 mixed with 60\%
incumbent raw predictions, then placed at the incumbent mean and spread.
The fixed candidate improved development ST-RAE by 0.008880 across 4/5
folds. On 220 held-out CYP2D6 labels it improved ST-RAE
0.831213→0.814041 and Spearman 0.54425→0.55725. Removing any one NB33
fold model preserved the gain. The paired scaffold-bootstrap 95\%
interval for the error difference is {[}-0.035570,+0.000300{]}, which
includes no improvement, and the holdout has been reused. The evidence
supported a small incremental candidate.

The rough forecast was 0.463, which was 0.4669 minus the local
four-target holdout improvement of 0.004293.

{\def\LTcaptype{none} % do not increment counter
\begin{longtable}[]{@{}
  >{\raggedright\arraybackslash}p{(\linewidth - 4\tabcolsep) * \real{0.3333}}
  >{\raggedright\arraybackslash}p{(\linewidth - 4\tabcolsep) * \real{0.3333}}
  >{\raggedright\arraybackslash}p{(\linewidth - 4\tabcolsep) * \real{0.3333}}@{}}
\toprule\noalign{}
\begin{minipage}[b]{\linewidth}\raggedright
Additional test
\end{minipage} & \begin{minipage}[b]{\linewidth}\raggedright
Measured result
\end{minipage} & \begin{minipage}[b]{\linewidth}\raggedright
Decision
\end{minipage} \\
\midrule\noalign{}
\endhead
\bottomrule\noalign{}
\endlastfoot
CYP2D6 potency-only and related-assay specialists, five folds each &
Best addition to v32 reduced endpoint development error only 0.000446,
3/5 folds; neither improves the fixed NB33 candidate & Exclude \\
NB29 control versus extra TDI labels, five paired small-model folds & C1
gain only 0.000704, 3/5 folds; C9/C3 best weight zero & No full-size run
justified by this pilot \\
NB33 C1 at 10\%, overlooked by best-before-consistency selection &
Development improved 5/5 folds; fixed holdout worsened 0.745284→0.746917
& Exclude \\
C1 saved-model complement, selected from 38 development comparisons &
Fixed holdout worsened 0.745284→0.752635; rho also declined & Exclude,
despite development gain \\
C1 potency-category prediction versus matched regression & Best
interval-action development gain 0.002542, only 2/5 folds & Exclude \\
\end{longtable}
}

The original NB33 breakout gate remains failed. The incremental checks
were separately frozen and authorized; they do not retroactively turn
that gate into a pass. The selection-order finding was a limitation of
the documented rule, not code contradicting the protocol. Local training
peaks were about 0.52GB for specialists and 0.80GB for the TDI pilot.

\textbf{Later records in this folder take precedence over earlier
pending plans.} The v34 upload (v33 with spread rescaling on CYP1A2 and
CYP3A4 and a CYP2C9 blend) scored \textbf{0.4685, rank 41}, worse than
v32, so v32 remains the best uploaded file. Its per-isoform errors were
0.5184 for CYP1A2, 0.3767 for CYP2C9, 0.5728 for CYP2D6 and 0.4061 for
CYP3A4. CLAMP embeddings were not evaluated.

\textbf{Tox21 CYP1A2 as an external predictor (AID 1671199).} The assay
holds 7,671 unique compounds, and the luciferase counter-screen (AID
1224835) flagged 864 rows, which were removed. A gradient-boosting model
trained on 4,995 of the remaining compounds reached Spearman 0.411 on
this project's CYP1A2 labels against 0.564 for the incumbent arm,
correlated 0.052 with that arm's residual, and moved holdout Spearman
from 0.5884 to 0.5905. Only 78 labelled challenge compounds overlap the
assay, and the two assays agree on them at Spearman 0.642.

\textbf{Placement inside the credible interval.} The scoring function
charges nothing for a prediction inside a compound's credible interval.
In the training data the weak CYP2D6 compounds (pIC50 below 3.6) have a
median upper bound of 3.75, while the CYP2D6 column in v34 has a median
of 3.20 with 65\% of its values below 3.5. An offline test used the
labelled development compounds with their real intervals plus blank
cells whose screen showed no inhibition, given an assumed window from
1.06 to a variable upper bound (64\% of that set is inactive, close to
the share expected in the blind set). Pulling values below 3.5 toward
3.5 while keeping their order moved the error from 0.557 to 0.517 when
the window top is 3.75 and from 0.591 to 0.550 when it is 3.52, and left
it at 0.643 when the window top is 3.3. The unmodified column scored
0.557 in this setup against 0.5728 on the board. The assumed values for
inactive compounds make this a simulation, and the file carrying the
change had not been scored when this was written.

\textbf{What public write-ups report.} Scigantic reports macro error
0.4436 at rank 11 on the interim board from twelve members, eight
CheMeleon-based and four Chemprop-based, combined with per-endpoint
weights fitted on out-of-fold rows and rescaled to the incumbent's mean
and spread
(\href{https://github.com/Scigantic/open-admet-challenge-cyp}{repository}).
The Workbench entry trains four stock Chemprop networks on 31,670
compounds from the challenge, ChEMBL and the Veith qHTS data and
recalibrates each isoform with an affine map
(\href{https://supercowpowers.github.io/workbench/blogs/cyp_challenge/}{write-up}).
On the live board of 27 October the leading per-isoform errors were
0.3659 for CYP1A2, 0.2776 for CYP2C9, 0.3952 for CYP2D6 and 0.3162 for
CYP3A4. This entry's CYP2D6 column matches entries scoring 0.47 to 0.48
on MAE, R-squared and Spearman and scores 0.573.

Detailed measurements, provenance and restart notes are in
\href{2026-10-26-breakout-continuation-round2.md}{the dated continuation
record}. The method narrative above is the 23 October snapshot; its
version-specific descriptions are not the current candidate recipe.

\subsection{Final week: 27 to 29
October}\label{final-week-27-to-29-october}

\textbf{Seeds of the NB33 recipe, under a frozen gate.} Two more
five-fold runs of NB33 (seed blocks 2042 and 3042) and a variant with
lighter auxiliary weights and lower dropout (NB36) were judged by one
rule fixed in advance: blend weights 0, 0.1, 0.2, 0.4 or 0.6 against the
incumbent, at least four of five folds improving, and a combined
development gain of at least 0.005. NB33 alone gained 0.0022 and failed.
NB33 averaged with both seeds gained 0.0062 and passed. Adding NB36
lowered that to 0.0057: alone it matched NB33, and in the average it
added nothing. The CYP1A2 blend from this route reached the board in
\texttt{v38} and lowered that column from 0.5184 to 0.4932, a larger
gain than development had forecast.

\textbf{Weak-compound placement, read off the board one column at a
time.} Development folds contain no labels for compounds that never
inhibited, so the only measurement of where to put them is the board.
Each isoform is scored on its own tab, so an upload that changes one
column measures that change alone. CYP2D6 predictions below 3.5 were
pulled toward 3.5 by the order-preserving map x → 3.5 − g(3.5 − x):

{\def\LTcaptype{none} % do not increment counter
\begin{longtable}[]{@{}lrrr@{}}
\toprule\noalign{}
Strength g & 0.5 & 0.3 & 0.15 \\
\midrule\noalign{}
\endhead
\bottomrule\noalign{}
\endlastfoot
CYP2D6 tab & \textbf{0.5524} & 0.5556 & 0.5714 \\
\end{longtable}
}

The best reading is at g = 0.5, and differences this small sit close to
the noise of a board that scores half of the blind set. \texttt{v39}
changed all four columns at once, adding a floor at 3.6 on the other
three isoforms and the 0.15 strength on CYP2D6, and every tab got worse
(CYP2D6 by 0.016), which cost a slot that taught little; after it, no
upload changed a column without evidence behind the change.

\textbf{A descriptor-augmented network was built and not run.} NB37 adds
216 RDKit descriptors to the NB33 network and passed a local smoke test.
Given that every earlier family had been null or close to it, the three
to five GPU hours were not spent.

\textbf{Choosing the last network by asking what the stack was missing.}
Rather than guess at a fifth architecture, each of Scigantic's thirteen
published member columns was added to this project's combiner one at a
time and scored on held-out folds. The columns that helped were the
randomly initialised and the precision-weighted ones, their arm B most
of all (0.019 on CYP2C9, 0.017 on CYP2D6), while their pretrained
members added nothing this stack did not already have. NB38 was built to
that specification: the NB33 network with the CheMeleon initialisation
removed, trained on absolute error across all twenty-two heads, with
each compound weighted by the mean inverse variance of its credible
intervals, normalised to the median and clipped to the range 0.25 to 4.
Thirty-seven per cent of the 4,179 labelled compounds reach a clip,
which is how wide the spread of interval widths is. Five folds took
about five GPU hours. In the combiner it took the largest weight on
CYP2C9 and CYP2D6 and the second largest on CYP1A2, and about a third of
the own-only board gain is attributable to it.

\textbf{CYP2D6 has now failed to transfer four times.} A pseudo-label
arm improved the locked holdout and the board disagreed; NB33's gate
gain of 0.0172 on the holdout came back 0.007 worse on the board; the
combiner predicted 0.5282 for the blend and the board returned 0.5449;
and the own-only column moved from 0.5524 to 0.5669 while development
said it would improve. Every other isoform has transferred in the same
direction as development throughout. The cause is the population rather
than the method: development can only score compounds that carry a
label, those compounds are the hits, and the blind half contains the
weak compounds whose ordering this project's CYP2D6 models get wrong.
The working rule at the end of this project is that a development gain
on CYP2D6 is not evidence, and the column is chosen by whichever version
the board has scored best.

\subsection{Validation, and where it misled
us}\label{validation-and-where-it-misled-us}

A locked holdout of 726 compounds, scaffold-disjoint from development,
was scored once per candidate. Across three submissions where the
holdout and the leaderboard can be compared, the offset between them was
−0.0834, −0.0728 and −0.0833, so two of three agreed to within 0.0001
and predictions from the third submission onward were accurate to
0.0004.

The exception matters. One change improved the holdout by 0.0046 with a
paired bootstrap putting the probability of improvement at 0.90, and
then scored 0.0060 worse on the blind set. The pattern that has held
since is that confidence intervals excluding zero transfer and intervals
crossing zero do not, which is a small sample and an uncomfortable rule
to rely on. A 726-compound holdout cannot resolve effects of 0.005, and
several decisions in this project were made as though it could.

A second exception runs the other way and is larger. The holdout
contains only hit compounds, and it scored the pseudo-labelled CYP2D6
arm 0.019 worse and the higher-weight blend a further 0.021 worse, while
the leaderboard improved that column by about 0.053 and then by about
0.025. For an isoform whose blind set differs from the labelled
population, the holdout is a sample of the wrong population, and its
sign cannot be trusted for changes aimed at that difference.

\subsection{Limitations}\label{limitations}

The stack built here, uploaded alone as \texttt{v46}, scored 0.4472 at
rank 45 of about 283 when it was read. That standing depends on Scigantic's
published predictions, and without them this entry sits outside the top
ten percent of the live board. Their own standing was 0.4283 at rank 21 when it was last
read on 28 October, from a submission later than the R25 file used here,
so the submitted file is ahead of the team it borrows from by 0.0091,
which is around four times the half-set noise. That margin was produced
by calibration applied after the blend rather than by the blend itself. Entries above \texttt{v46} run six base model families
against this one's three, pretrain every backbone on public bioactivity
at scale, and in one case train on tens of thousands of unlabelled
catalogue molecules chosen to resemble the blind set. None of those is a
trick; they are infrastructure accumulated over a longer period.

The findings above are measured on one challenge with 4,905 training
compounds and one blind set of 750, and the ones that were confirmed on
blind data are confirmed by a single submission each.

\subsection{Reproduction}\label{reproduction}

Code and the final submission pipeline are in this repository. The
supporting record, including every closed direction with its measurement
and the errors made along the way, is in \texttt{STATUS.md} and
\texttt{METHOD.md}.

\end{document}
